
\section{Equipartitioning}
\label{sec:6.5}
% BEGIN SUBJECT INDEX: sec:6.5
\index{equipartition!kernel blocks}
\index{equipartition!stationary measure}
% END SUBJECT INDEX


We show basic results based on the fact due to BM-F 
that $\pi k=(1/r)\rho(k)$, for every $k\in\cK$. Given this, it follows that, with $e\in \E(\cK)$ such that
$ek=ke=k$, $\pi kk^*=\pi ee^*=(1/r)u$, and averaging yields 
$$
  \pi N = (1/r) u
$$
Now suppose that $f$ is any idempotent of $\cK$ such that $\rho(k)=\rho(f)$. Then 
$$
  \pi kf^*=(1/r)\,\rho(k) f^*=(1/r)\,
  \rho(f) f^*=\pi ff^*=(1/r)u
$$ 
Taking $k=we$, we have the useful form
$$
  \pi w ee^*=\frac{1}{r}\,u
$$
for any $w\in\cS$, and any $e\in \E(\cK)$, since $we\in\cK$ has the same range as $e$.

We now show these results directly.

\begin{proposition} 	\label{prop:6.2}
For every $e\in\E(\cK)$, $\pi\Delta e=0$. Equivalently, $\pi Re=\pi Be=\pi e$.
\end{proposition}

\begin{proof} Multiplying eq.\tsp(\ref{eq:N}), for $N$,
 on the left by $\Omega$ and on the right by $\Omega^*$, we have,
using $\Omega A=A\Omega=\Omega$,
$$
  \sym{\Omega}{N}+\Omega\sym{\Delta}{N}\Omega^*
	=\sym{\Omega}{N}
$$
Hence,
$$
  \Omega\sym{\Delta}{N}\Omega^*=0
	=\alpha_\mu\Omega\Delta e_\mu e_\mu^*\Delta^*\Omega^*
$$
Taking traces, all terms are nonnegative so that $\Omega\Delta e=0$ for every $e$. From the definition of $\Delta$ the alternative
formulation results after averaging $\pi Re=\pi Be$ to $\pi Ae=\pi e$.  \qed
\end{proof}\label{proppidelta}

\begin{corollary} 	\label{cor:6.1}\label{ch5:kernel-delta}
For any $k\in\cK$, $\pi\Delta k =0$.
\end{corollary}

This follows from Proposition~\ref{prop:6.2} by writing $k=ek$, where $e$ is the identity of the local group containing $k$, so $ek=ke=k$. As well

\begin{corollary} \label{cor:pideltan}  \label{cor:6.2}
We have $\pi\Delta \bigl(J-N\bigr) =0$.
\end{corollary}

\begin{proposition} 	\label{prop:6.3}
For every $w\in\cS$, and $e\in\E(\cK)$, $\pi we=\pi e$.
\end{proposition}

\begin{proof} This follows by induction on $l$, the length of $w$. Proposition \ref{proppidelta}
is the result for $l=1$. Let $w_{l+1}$ be a word of
length $l+1$, then $w_{l+1}=$ one of $Rw_l$, $Bw_l$, where $w_l$ has length $l$. Since $w_le\in\cK$ for $e\in\cK$, Corollary~\ref{ch5:kernel-delta} yields 
\begin{eqnarray*}
\pi w_{l+1} e &=&\pi R w_le\cr &=& \pi B w_le \cr &=& \pi w_l e \hbox to 100pt{}{\rm (averaging)}\cr
&=& \pi e \hbox to 100pt{} \hbox{\rm (by induction)}\qed\cr
\end{eqnarray*}
\end{proof} \label{proppwe}

\begin{proposition}  	\label{prop:6.4}
For any $e\in\E(\cK)$, $w\in \cS$, 
$$
  \pi we=(1/r)\,\rho(e) \qquad \hbox{\ and\ }
	\qquad\pi w ee^*=(1/r)\,u
$$
%
\end{proposition}

\begin{proof} From Proposition \ref{proppidelta}, 
multiplying on the right by $e^*$ , we have $\pi\Delta ee^* =0$. 
Averaging over $e$ yields $\pi\Delta N = 0$. Equivalently, $N\Delta^*\dg\pi.   = 0$. 
Hence, multiplying equation \eqref{eq:N} on the right by  $\dg\pi.$ , 
using $A^*\dg\pi. = \dg\pi.$, we find
$$
   AN\dg\pi.  = N \dg\pi.
$$
So $N\dg\pi.$  is fixed by $A$, hence is a constant vector. Transposing back, say
$\pi N = au$ for some constant $a$. Now let $e\in \E(\cK)$ and take the $e_i$ with the
same range as $e$ in the summation defining $N$. Then from Proposition \ref{proppwe} we have
$\pi e e_i=\pi e_i$. Since the idempotents of a column are a left-zero semigroup, $\pi ee_i=\pi e$. Thus
on the one hand $ \pi eNe^*=\pi N e^*$
and on the other $\pi eNe^*=\pi ee^*$.
Or $\pi ee^* =\pi Ne^*=aue^*=au$ for any idempotent $e$.  Thus, $\pi e=a\rho(e)$ and
$\pi e\dg u.=1=ar$ yields $a=1/r$, where $r$ is the rank of the kernel. So, for any $w\in\cS$,
$\pi we=\pi e = (1/r)\rho(e)$ and $\pi wee^*=(1/r)u$, as required. \qed
\end{proof}

Averaging the second relation in the above Proposition

\begin{corollary} \label{cor:pideltann}  \label{cor:6.3}
For any $w\in\cS$,
$\pi w N =\frac{1}{r}\,u$. And $\pi w \bigl(J-N\bigr) =r'\,u$.
\end{corollary}
